\documentclass[12pt,a4paper]{article}
% Компилируйте в Overleaf с XeLaTeX.
\usepackage{fontspec}
\setmainfont{Noto Serif}
\usepackage[russian]{babel}
\usepackage{geometry}
\usepackage{amsmath,amssymb,mathtools}
\usepackage{array,booktabs}
\usepackage{xcolor}
\usepackage{tikz}
\usepackage[hidelinks]{hyperref}

\geometry{left=20mm,right=20mm,top=22mm,bottom=22mm}
\setlength{\parindent}{0pt}
\setlength{\parskip}{0.55em}
\renewcommand{\arraystretch}{1.4}
\sloppy

\begin{document}

% Титульный лист
\begin{titlepage}
\centering
\vspace*{0.29\textheight}
{\LARGE\bfseries Ревью домашней работы\par}
\vspace{2.3cm}
{\large Автор: Лёвин Артём Александрович\par}
\vspace{0.8cm}
{\large 9 октября 2026 года\par}
\vfill
\begin{tikzpicture}
  \draw[blue!30, line width=0.8pt]
    (-3,0) .. controls (-1.6,0.35) and (-0.7,-0.35) .. (0,0)
           .. controls (0.7,0.35) and (1.6,-0.35) .. (3,0);
\end{tikzpicture}
\vspace*{1.1cm}
\end{titlepage}

% Сводная таблица: только задания с ошибками
\newpage
\begin{center}
{\Large\bfseries Сводная таблица ошибок}
\end{center}
\vspace{1cm}
\small
\begin{tabular}{@{}p{0.08\linewidth}p{0.17\linewidth}p{0.27\linewidth}p{0.16\linewidth}p{0.16\linewidth}@{}}
\toprule
\textbf{№} & \textbf{Тема} & \textbf{Суть ошибки} &
\textbf{Ваш ответ} & \textbf{Правильный ответ}\\
\midrule
\hyperref[task:6]{6} &
Свойства степеней &
При делении степеней не учтён отрицательный показатель степени тройки. &
\(6\) & \(\dfrac{2}{3}\)\\
\bottomrule
\end{tabular}
\normalsize

% Разделительная страница
\newpage
\thispagestyle{empty}
\null\vfill
\begin{center}
{\Huge\bfseries Разбор ошибок}
\end{center}
\vfill\null

% Подробный разбор
\newpage
\phantomsection\label{task:6}
\section*{Задание № 6}

\textbf{Текст задания.}

Упростите выражение и найдите его значение:
\[
\frac{2^{7/2}\cdot 3^{3/2}}{6^{5/2}}.
\]

\textbf{Ваше решение.}

Вы верно представили шестёрку как произведение двойки и тройки
и применили правило степени произведения:
\[
\frac{2^{7/2}\cdot3^{3/2}}{6^{5/2}}
=
\frac{2^{7/2}\cdot3^{3/2}}{(3\cdot2)^{5/2}}
=
\frac{2^{7/2}\cdot3^{3/2}}
     {3^{5/2}\cdot2^{5/2}}.
\]
После этого в записи сделан переход к выражению
\[
2\cdot3=6.
\]

\textbf{Ваш ответ.} \(\boxed{6}\).

\textcolor{red}{\textbf{Ошибка:}} при сокращении степеней
с основанием \(3\) получился множитель \(3\), хотя
должен получиться множитель \(\dfrac13\).

\textbf{Где именно возникает ошибка.}

Последний правильный шаг:
\[
\frac{2^{7/2}\cdot3^{3/2}}
     {3^{5/2}\cdot2^{5/2}}.
\]
Первый неправильный переход --- замена этой дроби
на произведение \(2\cdot3\).
Деление степеней с одинаковым основанием требует
\emph{вычесть показатель знаменателя из показателя числителя}.
Для тройки получаем:
\[
\frac{3^{3/2}}{3^{5/2}}
=3^{\,3/2-5/2}
=3^{-1}
=\frac13,
\]
а вовсе не \(3\).

\textbf{Почему этот шаг неверен.}

По правилу деления степеней при положительном основании \(a\)
\[
\frac{a^m}{a^n}=a^{m-n}.
\]
Это правило действует и тогда, когда разность показателей
отрицательна. Отрицательный показатель не означает,
что степень становится отрицательным числом.
Он указывает на обратную величину:
\[
a^{-k}=\frac{1}{a^k}\quad (a\ne0).
\]
В рассматриваемой дроби в числителе у тройки показатель
равен трём вторым, а в знаменателе --- пяти вторым.
Следовательно, после деления показатель равен
минус единице. Поэтому остаётся не три, а одна треть.
Если заменить одну треть на три, значение выражения
увеличится в девять раз: это и произошло в Вашем ответе.

\textcolor{blue!60!black}{\textbf{Ключевая идея:}}
при сокращении степеней отдельно работайте
с каждым основанием. Даже если показатель
после вычитания оказался отрицательным,
правило деления не меняется.

\textbf{Исправление от последнего правильного шага.}

Сохраним Ваше верное разложение знаменателя
и сгруппируем множители с одинаковыми основаниями:
\begin{align*}
\frac{2^{7/2}\cdot3^{3/2}}
     {3^{5/2}\cdot2^{5/2}}
&=\frac{2^{7/2}}{2^{5/2}}
  \cdot\frac{3^{3/2}}{3^{5/2}}\\
&=2^{\,7/2-5/2}\cdot3^{\,3/2-5/2}\\
&=2^1\cdot3^{-1}\\
&=2\cdot\frac13=\frac23.
\end{align*}

\textbf{Полное правильное решение.}

Сначала заметим, что основание степени в знаменателе
составное: \(6=2\cdot3\). По свойству степени произведения
\[
6^{5/2}=(2\cdot3)^{5/2}
       =2^{5/2}\cdot3^{5/2}.
\]
Подставим это в исходное выражение:
\[
\frac{2^{7/2}\cdot3^{3/2}}{6^{5/2}}
=
\frac{2^{7/2}\cdot3^{3/2}}
     {2^{5/2}\cdot3^{5/2}}.
\]
Вычтем показатели степеней двойки:
\[
\frac72-\frac52=\frac{7-5}{2}=1.
\]
Затем вычтем показатели степеней тройки:
\[
\frac32-\frac52=\frac{3-5}{2}=-1.
\]
Значит,
\[
2^{\,1}\cdot3^{-1}
=2\cdot\frac{1}{3}
=\boxed{\frac23}.
\]

\textbf{Проверка.}

Умножим полученный ответ на исходный знаменатель:
\begin{align*}
\frac23\cdot6^{5/2}
&=\frac23\cdot2^{5/2}\cdot3^{5/2}\\
&=2^{\,1+5/2}\cdot3^{\,-1+5/2}\\
&=2^{7/2}\cdot3^{3/2}.
\end{align*}
Получили в точности числитель исходной дроби.
Следовательно, вычисление верно.

\textcolor{teal!60!black}{\textbf{Итог:}}
правильное значение выражения --- \(\boxed{\dfrac23}\).
Ошибка возникла не при разложении шестёрки,
а при применении правила деления степеней тройки.

\subsection*{Задача на закрепление}
Упростите выражение и найдите его значение:
\[
\frac{5^{7/2}\cdot2^{3/2}}{10^{5/2}}.
\]

\vfill
\begin{center}
\small Лёвин Артём Александрович эксклюзивно для Анны Банновой
\end{center}

\end{document}
